% s-170-p170-m002; original PDF 170, 171, 172
\Aphorism{s-170-p170-m002}{258}
\Sperr{Auf der Strasse.} Ich bin euer und ihr seid mein; was thut’s, dass wir einander nicht kennen? Sich kennen heisst sich schon zu nahe getreten sein; es ist die gröbste Form von Beziehung. Wir halten gerade die Mitte zwischen Fremd- und Bekanntsein; wir gehen ohne Gruss auf der Strasse an einander vorüber, und doch ist uns etwas gemeinsam, eben die Strasse, die Atmosphäre, der grosse und kleine Habitus irgend welcher Stadt, die gerade gross und klein genug ist, um etwas „Welt“ mit etwas „Heimath“ zu vereinigen. Diese vielen Gesichter, denen man gewohnheitsmässig begegnet und die sich durch Wiederholung im Gehirn einzeichnen, haben darum noch keinen Nexus zu meinem Fühlen hinüber; aber eine enge Auswahl darunter spinnt auch diesen Faden an und so bilden wir eine kleine Gesellschaft, wo Keiner des Anderen Namen weiss. Wir würden den Blick vermissen, den wir zu wechseln pflegen, worin Jeder dem Anderen die Geschichte seines letzten Tages erzählt und sich ihm zur stillschweigenden Kritik hingiebt; ich würde den Mädchenkopf vermissen, der sich bei meinem Herankommen regelmässig abwendet, mit einem Ausdruck weihevoller Jungfräulichkeit, der mir regelmässig ein Lächeln abzwingt. Manche dieser Gesichter wollen über die Strassenbekanntschaft hinaus, sie machen Miene zu beichten, zu fragen, sich interessirt und sympathisch zu geben; andere wieder weisen ab und üben ihr deutsches Recht auf herausfordernde Unart — ich hüte mich wohl, das schwebende Verhältniss durch ein Mehr oder Weniger zerstören zu lassen. Ich rathe nicht an Romanen herum, von denen ich nur ein Endglied, den Gesichtsausdruck eines Betroffenen, zur Verfügung habe; ich folge Keinem aus der bewegten Strassenluft, durch die er mit mir communicirt, in die stille dumpfe Atmosphäre seines Eigenen, in der ich sofort zum Fremden würde. Der Gedanke, dass ich mich mit diesen Vielen, Unbekannten in ein Hectoliter Luft, in ein Netzhautbild, in die feinen unspürbaren Schwingungen des Milieus zu theilen habe, vermag mich für sie zu einer zarten Sympathie zu stimmen. Damit wäre es sofort vorbei, wenn die Gemeinsamkeit eine gröbere würde. Es besteht eine Fernwirkung zwischen uns, die bei der Berührung erlischt, wie zwischen elektrischen Ladungen, die sich bei der Berührung sofort neutralisiren. Ich könnte mich dieser seltsamen Genossenschaft, die ich als Ganzes, als körperliche Einheit empfinde, hingeben, mich für sie opfern; frage ich mich aber, wie und zu welchem Zwecke, zerfälle ich jenes Collectivum wieder in seine Individuen, die mich persönlich gar nichts angehen, so lache ich über sie und mich. Ich fühle mich in geheimer Weise mit ihnen mitschuldig, für sie mitverantwortlich; wären sie eine wirkliche Gruppe, zu der mich Interesse oder Gesinnung zöge, so würde ich gegen sie meinen Egoismus kehren und mich kühl verhalten. Uns Alle zusammen vermag ich mir als einen Bund Verschworener zu denken, die naturgemäss zu einander gehören; aber ich wüsste keinen Gedanken, kein Leiden, kein körperliches oder seelisches Abzeichen zu nennen, an dem wir uns erkennen sollten, keine Fremde und Feindschaft, gegen die wir uns als Landsleute und Gleichgesinnte zu betonen hätten. Ich bin der Eure und ihr seid die Meinen: ich machte euch eben eine Liebeserklärung, die ich sofort zurückziehe, wenn sich von einem Einzelnen unter euch zu mir etwas Persönliches anknüpft. Ich tausche mit euch strahlende Wärme aus; aber kalt wie Eis würde ich werden, wenn sich von euch herüber eine Menschenhand nach meiner streckte \Dots{3}
\par

